\documentclass[a4paper,11pt]{article}
\usepackage{amssymb, enumitem}
\setlist[enumerate,1]{label={\bfseries\arabic*.},ref={\bfseries{Ejercicio \arabic*}}}
\setlist[enumerate,2]{label={\bfseries(\alph*)},ref={\bfseries(\alph*)}}


\parindent 0cm
\usepackage{amssymb,amsmath,amsthm,latexsym,epsfig,euscript,multicol}
\usepackage{graphbox}

\usepackage[utf8x]{inputenc}
\usepackage{listings,xcolor,bm}


\definecolor{mygreen}{rgb}{0,0.6,0}
\definecolor{mygray}{rgb}{0.5,0.5,0.5}
\definecolor{mymauve}{rgb}{0.58,0,0.82}
\lstset{
  backgroundcolor=\color{white},   % choose the background color; you must add
  basicstyle=\small\ttfamily,      % the size of the fonts that are used for the code
  breakatwhitespace=false,         % sets if automatic breaks should only happen at whitespace
  breaklines=true,                 % sets automatic line breaking
  captionpos=b,                    % sets the caption-position to bottom
  commentstyle=\color{mygreen},    % comment style
  deletekeywords={...},            % if you want to delete keywords from the given language
  escapeinside={\%*}{*)},          % if you want to add LaTeX within your code
  extendedchars=true,              % lets you use non-ASCII characters; for 8-bits encodings only, does not work with UTF-8
  firstnumber=1,                % start line enumeration with line 1000
  frame=none,	                   % adds a frame around the code
  keepspaces=true,                 % keeps spaces in text, useful for keeping indentation of code (possibly needs columns=flexible)
  keywordstyle=\color{blue},       % keyword style
  language=Python,                 % the language of the code
  morekeywords={*,...},            % if you want to add more keywords to the set
  numbers=none,                    % where to put the line-numbers; possible values are (none, left, right)
  numbersep=5pt,                   % how far the line-numbers are from the code
  numberstyle=\tiny\color{mygray}, % the style that is used for the line-numbers
  rulecolor=\color{black},         % if not set, the frame-color may be changed on line-breaks within not-black text (e.g. comments (green here))
  showspaces=false,                % show spaces everywhere adding particular underscores; it overrides 'showstringspaces'
  showstringspaces=false,          % underline spaces within strings only
  showtabs=false,                  % show tabs within strings adding particular underscores
  stepnumber=5,                    % the step between two line-numbers. If it's 1, each line will be numbered
  stringstyle=\color{mymauve},     % string literal style
  tabsize=4,	                   % sets default tabsize to 2 spaces
  title=\lstname                   % show the filename of files included with \lstinputlisting; also try caption instead of title
}
% Caracteres especiales
\def\A{\mathbb{A}}
\def\C{\mathbb{C}}
\def \N{\mathbb{N}}
\def \P{\mathbb{P}}
\def \Q{\mathbb{Q}}
\def \R{\mathbb{R}}
\def \Z{\mathbb{Z}}
\def \sen{\textrm{sen}}

\def\Np{$\N$}
\def\Zp{$\Z$}
\def\Qp{$\Q$}
\def\Rp{$\R$}
\def\Cp{$\C$}

\def\bb{\bm{b}}
\def\bu{\bm{u}}
\def\bv{\bm{v}}
\def\bx{\bm{x}}
\def\bA{\bm{A}}
\def\bB{\bm{B}}
\def\bD{\bm{D}}
\def\bE{\bm{E}}
\def\bM{\bm{M}}
\def\bT{\bm{T}}


\def\K{\textrm{K}}
\def\V{\textrm{V}}
\def\S{\textrm{S}}

\def\degres{$^\circ$}

\newcount\todno
\def\no{\global\advance\todno by 1 \the\todno}

\topmargin-2cm \vsize 29.5cm \hsize 21cm
\setlength{\textwidth}{16.75cm}\setlength{\textheight}{23.5cm}
\setlength{\oddsidemargin}{0.0cm}
\setlength{\evensidemargin}{0.0cm}


\theoremstyle{definition}
\newtheorem{ejer}{Ejercicio}
\newcommand{\bej}{\begin{ejer}}
\newcommand{\fej}{\end{ejer}}

\begin{document}

\centerline{{\small Universidad de Buenos Aires - Facultad de Ciencias Exactas y Naturales - Ciencias de Datos}}

\vskip 0.2cm

\hrule

\vskip 0.2cm

 \centerline{{\bf\Large{\sc Laboratorio de Datos}}}

 \vskip 0.2cm

 \centerline{\ttfamily Primer Cuatrimestre 2024}

\vskip 0.2cm

 \hrule

 \bigskip
 \centerline{\bf Práctica N$^\circ$ 9: Descenso por gradiente y Redes Neuronales}
 \bigskip


Para realizar esta guía de ejercicios, descargar de la página de la materia el archivo \verb|tf_regressor.py| para poder correr el siguiente comando
\begin{lstlisting}
from tf_regressor import Regressor, train_test_split_scale_center
\end{lstlisting}

\begin{enumerate}
    \item En este ejercicio utilizaremos el dataset \verb|casos_coronavirus|.
    \begin{enumerate}
        \item Cargar el dataset y añadirle la columna \verb|dias_transcurridos| con el indice de cada observación
        \item Plotear \verb|dias_transcurridos| vs. \verb|confirmados_Nuevos|
        \item Queremos armar un modelo de regresión que permita explicar la evolución de casos de coronavirus ($Y$) en función de los días transcurridos ($X$). Para esto se proponen los siguientes modelos:
        \begin{enumerate}
            \item $Y = b + w_0\;X + w_1\; X^2$
            \item $Y = b + w_0\; X^{w_1}$ (tomar como valores iniciales $w_0=1,\, w_1=2,\, b$ aleatorio)
            \item $Y = b + w\; e^{X}$ 
            \item $Y = b + w_0\; e^{w_1\; X}$ (tomar como valores iniciales $w_0=1,\, w_1=1,\, b$ aleatorio)
        \end{enumerate}
        Dividir el conjunto de datos en entrenamiento y testeo y decidir qué modelo resulta más adecuado. Utilizar \verb|scikit-learn| para los modelos lineales y \verb|Regressor| para los no lineales. En este último caso,
        probar con distintas cantidades de épocas y otros valores iniciales de los pesos y del
bias.
        
        \textbf{Obs:} para el modelo iv) escribir \verb|f| utilizando \verb|np.e**(w[1]*x)| para $e^{w_1\; x}$
    \end{enumerate}


    \item En este ejercicio trabajaremos con el dataset \verb|titanic| de \verb|seaborn|
\begin{lstlisting}
titanic = sns.load_dataset('titanic')
\end{lstlisting}
\begin{enumerate}
\item Limpiando el dataset:
\begin{enumerate}
    \item contar cuantos \verb|NaN| tiene cada columna, y en base a eso decidir qué columna del dataset descartar antes de ejecutar \verb|.dropna()|
    \item graficar un boxplot de \verb|fare| (precio del boleto), ¿qué se observa?
    \item explorar el método de pandas \verb|quantile| para calcular el cuantil $0.99$ de la columna \verb|fare| y utilizarlo para eliminar las observaciones con outliers en esa columna.
\end{enumerate}
\item Realizar regresión logística para predecir la variable binaria de supervivencia (\verb|survived|) a partir del precio del boleto (\verb|fare|). ¿Qué porcentaje de casos clasifica correctamente?
\item Repetir el item anterior, considerando la interacción de la suma de \verb|fare| y \verb|age| con \verb|adult_male|. ¿Cuánto mejoró la precisión de la clasificación? ¿Qué se puede concluir a partir de la mejora en la precisión y del análisis de los pesos que el modelo otorga a cada variable?
\item Proponer un método que permita obtener una clasificación más precisa mediante regresión logística. Las demás columnas del DataFrame son:
\begin{itemize}
    \item \verb|pcclass| : clase en la que viajaba
    \item \verb|sibsp| : si viajaba con hermanos/as o cónyugues
    \item \verb|parch| : cantidad de hijes o padres con los que viajaba
    \item \verb|embarked| : donde se embarcó
    \item \verb|class| : nombre de la clase en la que viajaba (dato de \verb|pcclass| en string)
    \item \verb|embark_town| : nombre del lugar donde embarcó
    \item \verb|alive| : si sobrevivió (mismo valor que \verb|survived| pero booleano, por lo tanto no usar para predecir)
    \item \verb|alone| : si viajaba solo/a (es \verb|True| si $sibsp=0$ y $parch=0$)
\end{itemize}
\end{enumerate}

\item Utilizando el dataset \verb|diabetes.csv|, se quiere desarrollar un modelo para predecir si una persona tiene o no diabetes en base a las características descriptas en el resto de las columnas.
\begin{enumerate}
    \item Entrenar un perceptrón simple que permita categorizar a una persona como diabética o no diabética utilizando la sigmoidea como función de activación. ¿Cuántos falsos negativos hay? ¿Cómo se intepretan los pesos de la red entrenada?
    \item Diseñar una red multicapa para la clasificación, utilizando la sigmoidea como función de activación. ¿Cómo se desempeña este modelo con respecto al anterior?
\end{enumerate}

\item Probar que un perceptrón simple con la identidad como función de activación es equivalente a un modelo de Regresión Lineal.



\end{enumerate}
\end{document}

